

\usepackage{minted} % code listings
\usepackage{hyperref} % hyperlinks
\hypersetup{
    colorlinks=true,
    linkcolor=blue,
    urlcolor=cyan,
    filecolor=magenta,
    citecolor=green,
    pdfborder={0 0 0}
}

\title{Memory-safe Programming: An Introduction to Rust}
\subtitle{Rust at the University of Bamberg}
\author{Lukas Beierlieb M.Sc.}
\date{Summer Term 2026}

\begin{document}

\frame[plain]{\titlepage}

\begin{frame}[fragile]
    \frametitle{Copyright / Acknowledgements}

\end{frame}

\begin{frame}{Why Rust?}
    This is a slide.
\end{frame}

\begin{frame}[fragile]
    \frametitle{Hello World!}
    \begin{minted}[linenos]{rust}
fn main() {
    println!("Hello world");
}
    \end{minted}

    \pause

    Already in this tiny example:
    \begin{itemize}
        \pause
        \item Function definition (covering very soon)
        \item Rust program entry point
        \pause
        \item Declarative Macro
        \pause
        \item Unsized types
        \item References
        \item Lifetimes
    \end{itemize}
\end{frame}

\begin{frame}[fragile]
    \frametitle{Primitive Data Types}
    \begin{itemize}
        \item unsigned integer: \texttt{u8}, \texttt{u16}, \texttt{u32}, \texttt{u64}
        \item signed integer: \texttt{i8}, \texttt{i16}, \texttt{i32}, \texttt{i64}
        \item integers with architecture-dependent size: \texttt{usize}, \texttt{isize} (e.g., size of 64 bits on a 64-bit CPU)
        \pause
        \item floating point number: \texttt{f32}, \texttt{f64}
        \item \texttt{bool} with the values \texttt{true} and \texttt{false}
        \item \texttt{char}, which is 32 bits in size to contain Unicode characters
    \end{itemize}
    \pause
    Examples on the following slides!
\end{frame}

\begin{frame}[fragile]
    \frametitle{Creating Bindings with \texttt{let}}
    \begin{itemize}
        \item Informally called ``variables''
        \item Bind a value to a name
    \end{itemize}
    \pause
    \begin{minted}[linenos]{rust}
fn main() {
    let letter: char = 'a';
    let emoji: char = '\u{1F980}'
}
    \end{minted}
    \pause
    The compiler can often perform type inference.
    \pause
    \begin{minted}[linenos]{rust}
fn main() {
    let letter = 'a';
    let emoji = '\u{1F980}'
}
    \end{minted}
\end{frame}

\begin{frame}[fragile]
    \frametitle{Mutating Bindings}
    \begin{minted}[linenos]{rust}
fn main() {
    let letter = 'a';
    letter = 'b';
}
    \end{minted}
    \pause
    Compiler error:
    \begin{minted}{text}
error[E0384]: cannot assign twice to immutable variable `letter`
 --> src/main.rs:3:5
  |
2 |     let letter = 'a';
  |         ------ first assignment to `letter`
3 |     letter = 'b';
  |     ^^^^^^^^^^^^ cannot assign twice to immutable variable
    \end{minted}
\end{frame}

\begin{frame}[fragile]
    \frametitle{Binding Shadowing}
    \begin{minted}[linenos]{rust}
fn main() {
    let letter = 'a';
    let letter = 'b';
}
    \end{minted}
    \pause
    A binding shadows previous bindings of the same name.
    \pause
    \begin{minted}[linenos]{rust}
fn main() {
    let value = 'a';
    let value = true;
}
    \end{minted}
\end{frame}

\begin{frame}[fragile]
    \frametitle{Floating-point Numbers}
    \begin{minted}[linenos]{rust}
fn main() {
    let pi: f64 = 3.14159;
    let g: f32 = 9.81;
    let three: f64 = 3.;
}
    \end{minted}
    \pause
    You can omit the zero following the decimal point but not the decimal point.
    \pause
    \begin{minted}[linenos]{rust}
fn main() {
    let three: f64 = 3;  // will not compile
}
    \end{minted}
\end{frame}

\begin{frame}[fragile]
    \frametitle{Floating-point Numbers}
    If type inference cannot distinguish between \texttt{f32} and \texttt{f64}, the compiler defaults to \texttt{f64}.
    \pause
    \begin{minted}[linenos]{rust}
fn main() {
    let pi = 3.14159;  // pi's type is f64
}
    \end{minted}
\end{frame}

\begin{frame}[fragile]
    \frametitle{Unsigned Integers}
    \begin{minted}[linenos]{rust}
fn main() {
    let max_u8: u8 = 255;
    let max_u16: u16 = 65535;
    let max_u32: u32 = 4294967295;
    let max_u64: u64 = 18446744073709551615;
    let max_usize: usize = 18446744073709551615;  // on 64-bit platforms
}
    \end{minted}
\end{frame}

\begin{frame}[fragile]
    \frametitle{Signed Integers}
    \begin{minted}[linenos]{rust}
fn main() {
    let min_i8: i8 = -128;
    let min_i16: i16 = -32768;
    let min_i32: i32 = -2147483648;
    let min_i64: i64 = -9223372036854775808;
    let min_isize: isize = -9223372036854775808;  // on 64-bit platforms
}
    \end{minted}
\end{frame}

\begin{frame}[fragile]
    \frametitle{Alternative Integer Representations}
    \begin{minted}[linenos]{rust}
fn main() {
    let decimal = 123;
    let hexadecimal = 0xff00;
    let octal = 0o1277;
    let binary = 0b1101011;
}
    \end{minted}
\end{frame}

\begin{frame}[fragile]
    \frametitle{Integer Type Inference}
    If type inference cannot distinguish between integer types, the compiler defaults to \texttt{i32}.
    \pause
    \begin{minted}[linenos]{rust}
fn main() {
    let integer = 3;  // integer's type is i32
}
    \end{minted}
\end{frame}

\begin{frame}[fragile]
    \frametitle{Printing Values}
    We accept the functionality of the \texttt{println!} macro without looking deeper.
    \pause
    \begin{minted}[linenos]{rust}
fn main() {
    let integer = 3;
    let pi = 3.14;

    println!("The values are {}, {}, {}, {}", integer, pi, false, 'a');
}
    \end{minted}
    \pause
    Output:
    \begin{minted}{text}
The values are 3, 3.14, false, a
    \end{minted}
\end{frame}

\begin{frame}[fragile]
    \frametitle{Printing Values}
    Bindings can be included in the format string directly.
    \begin{minted}[linenos]{rust}
fn main() {
    let integer = 3;
    let pi = 3.14;

    println!("The values are {integer}, {pi}, {}, {}", false, 'a');
}
    \end{minted}
\end{frame}

\begin{frame}[fragile]
    \frametitle{Arithmetic Operators}
    \begin{minted}[linenos]{rust}
fn main() {
    let add = 3 + 4;
    let subtract = 10 - 5;
    let multiply = 2 * 7;
    let divide = 20 / 5;
}
    \end{minted}
\end{frame}

\begin{frame}[fragile]
    \frametitle{Logical Operators}
    \begin{minted}[linenos]{rust}
fn main() {
    let and = true && false;
    let or = false || true;
    let not = !true;
}
    \end{minted}
    \pause
    \texttt{\&\&} and \texttt{||} use short-circuit evaluation.
\end{frame}

\begin{frame}[fragile]
    \frametitle{Comparison Operators}
    \begin{minted}[linenos]{rust}
fn main() {
    let equal = 3 == 4;
    let not_equal = 5 != 7;
    let less_than = 10 < 100;
    let greater_than = 12 > 1;
    let less_or_equal = 11 <= 11;
    let greater_or_equal = 5 >= 30;
}
    \end{minted}
\end{frame}

\begin{frame}[fragile]
    \frametitle{Bitwise Operators}
    \begin{minted}[linenos]{rust}
fn main() {
    let and = 0xff & 0x7f00;
    let or = 0xffff | 0x100;
    let xor = 0x77 ^ 0x94;
    let not = !0xf;
    let shift_left = 0b101 << 2;
    let shift_right = 0b1000 >> 3;
}
    \end{minted}
\end{frame}

\begin{frame}[fragile]
    \frametitle{Blocks}
    Blocks are not regularly used on their own, but they are crucial for many other language constructs.

    \pause
    \begin{minted}[linenos]{rust}
fn main() {
    let x = 3;
    {
        let x = {};
    }
    println!("{x}");
}
    \end{minted}
    All \texttt{\{\ldots\}} (except \texttt{\{x\}} inside the string) are blocks.
\end{frame}

\begin{frame}[fragile]
    \frametitle{Blocks}
    \begin{minted}[linenos]{rust}
fn main() {
    let x = 3;
    {
        let x = 5;
        println!("x = {x}");
    }
    println!("x = {x}");
}
    \end{minted}
    \pause
    Output:
    \begin{minted}{text}
x = 5
x = 3
    \end{minted}
    \pause
    \begin{itemize}
        \item Bindings go out of scope at the end of the block they are declared in.
        \item Shadowed bindings can become available again.
    \end{itemize}
\end{frame}

\begin{frame}[fragile]
    \frametitle{Blocks}
    \begin{minted}[linenos]{rust}
fn main() {
    let x = {
        println!("print in the block");
        3 + 5
    };
    println!("x = {x}");
}
    \end{minted}
    \pause
    Output:
    \begin{minted}{text}
print in the block
x = 8
    \end{minted}
    \pause
    If a block ends with an expression, it evaluates to this expression.
\end{frame}

\begin{frame}[fragile]
    \frametitle{Blocks}
    \begin{minted}[linenos]{rust}
fn main() {
    let x = {
        let x = 3;
    };
    println!("x = {x}");
}
    \end{minted}
    \pause
    Compiler error:
    \begin{minted}{text}
error[E0277]: `()` doesn't implement `std::fmt::Display`
 --> src/main.rs:5:19
  |
5 |     println!("x = {x}");
  |                   ^^^ `()` cannot be formatted with...
    \end{minted}
\end{frame}

\begin{frame}[fragile]
    \frametitle{Blocks}
    \begin{minted}[linenos]{rust}
fn main() {
    let x = {
        let x = 3;
    };
    println!("x = {x}");
}
    \end{minted}
    \pause
    \begin{itemize}
        \item Line 3 is a statement, not an expression.
        \item The inner block defaults to returning to \texttt{()}.
        \item \texttt{println!} cannot directly print this value.
    \end{itemize}
\end{frame}


\begin{frame}[fragile]
    \frametitle{The \textit{unit} Type}
    The \textit{unit} type \texttt{()} has a single value \texttt{()}.
    \begin{minted}[linenos]{rust}
fn main() {
    let unit: () = ();
    let unit = ();  // type inference also works for this type
}
    \end{minted}
    \pause
    \begin{itemize}
        \item The type occurs often implicitly.
        \item Explicit uses of \texttt{()} exist.
    \end{itemize}
\end{frame}

\begin{frame}[fragile]
    \frametitle{Blocks}
    More examples for block evaluation:
    \begin{minted}[linenos]{rust}
fn main() {
    let a: () = {};  // empty block -> no expression -> evaluates to ()
    let b: () = {
        let x = 3;
    };  // let binding is a statement -> no expr. -> eval to ()
    let c: () = {
        3;
    };  // expression followed by semicolon is also a statement -> ()
    let d: usize = {
        3
    };  // block contains only expression -> eval to 3
    ...
    \end{minted}
\end{frame}

\begin{frame}[fragile]
    \frametitle{Blocks}
    \begin{minted}[linenos,firstnumber=11]{rust}
    ...
    let e: usize = {
        let x = 3;
        println!("x = {x}");
        x
    };  // block has statements and an expr. at end -> eval to 3
}
    \end{minted}
\end{frame}

\begin{frame}[fragile]
    \frametitle{Blocks}
    We have observed block examples such as:
    \begin{minted}[linenos]{rust}
fn main() {
    {
        3
    }
    let a = {};
}
    \end{minted}
    \pause
    When is a \texttt{;} required after a block?
\end{frame}

\begin{frame}[fragile]
    \frametitle{Blocks}
    \texttt{let} statements always have to end with a semicolon.
    \begin{minted}[linenos]{rust}
fn main() {
    let a = {};
    let b = {
        let x = 3;
        let y = 8;
        x + y
    };
}
    \end{minted}
\end{frame}

\begin{frame}[fragile]
    \frametitle{Blocks}
    Usually, expressions cannot exist between statements.
    \pause
    \begin{minted}[linenos]{rust}
fn main() {
    let a = 3;
    a;  // valid: expression with semicolon is a statement
    a   // invalid: standalone expression followed by something else
    let b = 3;
}
    \end{minted}
    \pause
    Probably for parsing reasons.
\end{frame}

\begin{frame}[fragile]
    \frametitle{Blocks}
    Blocks are also expressions, but they do not require a semicolon.
    \pause
    \begin{minted}[linenos]{rust}
fn main() {
    let a = 3;
    { a }  // block is an expression and evaluates to 3 but needs no ;
    { a };  // semicolon is allowed
    { a };;;  // actually, you can always place additional semicolons
    let b = 3;
}
    \end{minted}
    \pause
    Probably for parsing reasons.
\end{frame}

\begin{frame}[fragile]
    \frametitle{\texttt{if} Expressions}
    \begin{minted}[linenos]{rust}
fn main() {
    let age: u8 = 23;

    if age >= 18 {
        println!("Adult!");
    } else if age >= 10 {
        println!("Teenager!");
    } else {
        println!("Child!");
    }
}
    \end{minted}
\end{frame}

\begin{frame}[fragile]
    \frametitle{\texttt{if} Expressions}
    \begin{itemize}
        \item A condition must be an expression evaluating to a \texttt{bool}.
        \item Parentheses around conditions are not necessary.
        \pause
        \vspace{.5cm}
        \item \texttt{if} is an expression---it evaluates to the block of the chosen branch.
    \end{itemize}
    \pause
    \begin{minted}[linenos]{rust}
fn main() {
    let a = 3;
    let b = 5;
    
    let max = if a > b {
        a
    } else {
        b
    }
}
    \end{minted}
\end{frame}

\begin{frame}[fragile]
    \frametitle{\texttt{if} Expressions}
    \begin{minted}[linenos]{rust}
fn main() {
    let a = if true {
        3
    } else {
        'a'
    };
    println!("{}", a);
}
    \end{minted}
    \pause
    Compiler error:
    \begin{minted}{text}
error[E0308]: `if` and `else` have incompatible types
    \end{minted}
\end{frame}

\begin{frame}[fragile]
    \frametitle{\texttt{if} Expressions}
    Full compiler error:
    \begin{minted}{text}
error[E0308]: `if` and `else` have incompatible types
 --> src/main.rs:5:9
  |
2 |       let a = if true {
  |  _____________-
3 | |         3
  | |         - expected because of this
4 | |     } else {
5 | |         'a'
  | |         ^^^ expected integer, found `char`
6 | |     };
  | |_____- `if` and `else` have incompatible types
    \end{minted}
\end{frame}

\begin{frame}[fragile]
    \frametitle{\texttt{if} Expressions}
    \begin{minted}[linenos]{rust}
fn main() {
    if true {
        3
    } else {
        'a'
    }
    println!("hello");
}
    \end{minted}
    \pause
    Same compiler error.
    \pause
    \begin{itemize}
        \item Even if the value is not used, its type must be fixed.
        \pause
        \item As with blocks, a semicolon is not necessary.
    \end{itemize}
\end{frame}

\begin{frame}[fragile]
    \frametitle{\texttt{if} Expressions}
    Rust does not have the ternary operator (\texttt{?:}), but because \texttt{if} is an expression, you can use it similarly.
    \pause
    \begin{minted}[linenos]{rust}
fn main() {
    println!("{}", if true { 1 } else { 0 });
}
    \end{minted}
\end{frame}

\begin{frame}[fragile]
    \frametitle{\texttt{if} Expressions}
    \begin{itemize}
        \item \texttt{if} without \texttt{else}: The block must evaluate to \texttt{()}.
        \pause
        \item The following expressions are equivalent.
    \end{itemize}
    \begin{minted}[linenos]{rust}
fn main() {
    if true { 3 }
    if true { 3 } else {}
    if true { 3 } else { () }
}
    \end{minted}
    \pause
    Compiler error (for Line 2):
    \begin{minted}{text}
2 |     if true { 3 }
  |     ----------^--
  |     |         |
  |     |         expected `()`, found integer
  |     `if` expr. without `else` arms expect their inner expr. to be `()`
    \end{minted}
\end{frame}

\begin{frame}[fragile]
    \frametitle{\texttt{loop} Expressions}
    \begin{minted}[linenos]{rust}
fn main() {
    let mut counter: u64 = 0;
    loop {
        println!("loop iteration {counter}");
        counter += 1;
    }
}
    \end{minted}
    \pause
    \texttt{loop} will not stop!
\end{frame}

\begin{frame}[fragile]
    \frametitle{\texttt{loop} Expressions}
    \begin{minted}[linenos]{rust}
fn main() {
    let mut counter: u64 = 0;
    loop {
        counter += 1;

        if counter == 100 {
            break;
        }
    }
}
    \end{minted}
    \pause
    \texttt{loop} will not stop!
\end{frame}

\begin{frame}[fragile]
    \frametitle{\texttt{loop} Expressions}
    \texttt{loop} is an expression!
    \pause
    \begin{minted}[linenos]{rust}
fn main() {
    let mut counter: u64 = 0;
    let val: () = loop {
        counter += 1;

        if counter == 100 {
            break;
        }
    };
}
    \end{minted}
\end{frame}

\begin{frame}[fragile]
    \frametitle{\texttt{loop} Expressions}
    \texttt{break;} is equivalent to \texttt{break ();}.
    \pause
    \begin{minted}[linenos]{rust}
fn main() {
    let mut counter: u64 = 0;
    let val: () = loop {
        counter += 1;

        if counter == 100 {
            break ();
        }
    };
}
    \end{minted}
\end{frame}

\begin{frame}[fragile]
    \frametitle{\texttt{loop} Expressions}
    The same way, \texttt{break} can return values of all types.
    \pause
    \begin{minted}[linenos]{rust}
fn main() {
    let mut counter: u64 = 0;
    let mut cumulative_sum = 0;
    let result = loop {
        cumulative_sum += counter;
        counter += 1;

        if counter == 100 {
            break cumulative_sum;
        }
    };
}
    \end{minted}
\end{frame}

\begin{frame}[fragile]
    \frametitle{\texttt{while} Expressions}
    The \texttt{loop} examples looked like \texttt{while} loops with extra steps.
    \pause
    \begin{minted}[linenos]{rust}
fn main() {
    let mut counter: u64 = 0;
    let mut cumulative_sum = 0;
    while counter < 100 {
        cumulative_sum += counter;
        counter += 1;
    }
}
    \end{minted}
\end{frame}

\begin{frame}[fragile]
    \frametitle{\texttt{while} Expressions}
    \texttt{break;} also works inside of \texttt{while}.
    \pause
    \begin{minted}[linenos]{rust}
fn main() {
    let mut counter: u64 = 0;
    while counter < 100 {
        counter += 1;

        if counter == 10 {
            break;   // return early
        }
    }
}
    \end{minted}
    \pause
    \begin{itemize}
        \item But \texttt{break} cannot return a value in \texttt{while}.
        \item \texttt{while} expression always return \texttt{()}.
    \end{itemize}
\end{frame}

\begin{frame}[fragile]
    \frametitle{Functions}
    We have seen a function definition in every example.
    \pause
    \begin{minted}[linenos]{rust}
fn main() {
    println!("Hello world");
}
    \end{minted}
    \pause
    This example demonstrates:
    \begin{itemize}
        \item \texttt{fn} keyword indicates the definition of a named function
        \item A function needs to specify to block.
    \end{itemize}
\end{frame}

\begin{frame}[fragile]
    \frametitle{Functions}
    \begin{minted}[linenos]{rust}
fn pi() -> f64 {
    println!("Returning pi");

    3.14159
}
    \end{minted}
    \pause
    \begin{itemize}
        \item Functions always return a value.
        \item You must state the type of the return value after \texttt{->}.
        \pause
        \item The function typically evaluates to value of the block.
        \item Early return with \texttt{return} is possible.
    \end{itemize}
\end{frame}

\begin{frame}[fragile]
    \frametitle{Functions}
    
    \begin{itemize}
        \item Functions \textbf{always} return a value.
        \item The following function definitions are equivalent.
    \end{itemize}
    \pause
    \begin{minted}[linenos]{rust}
fn main() {}

fn main() -> () {}

fn main() -> () {
    ()
}
    \end{minted}
\end{frame}

\begin{frame}[fragile]
    \frametitle{Functions}
    
    \begin{minted}[linenos]{rust}
fn do_nothing() {}
fn main() {
    let a = do_nothing();
}
    \end{minted}
    \pause
    is equivalent to:
    \begin{minted}[linenos]{rust}
fn do_nothing() {}
fn main() {
    let a: () = do_nothing();
}
    \end{minted}
\end{frame}

\begin{frame}[fragile]
    \frametitle{Functions}
    \begin{minted}[linenos]{rust}
fn first_or_second(return_fst: bool, fst: usize, snd: usize) -> usize {
    if return_fst {
        fst
    } else {
        snd
    }
}
    \end{minted}
    \pause
    \begin{itemize}
        \item Functions specify zero or more arguments.
        \item Type declarations of arguments are mandatory.
        \item Each argument is a binding available in the function's block.
    \end{itemize}
\end{frame}

\begin{frame}[fragile]
    \frametitle{Functions}
    \begin{minted}[linenos]{rust}
fn some_function(a: usize) -> usize {
    a = 5;

    a
}
    \end{minted}
    \pause
    Compiler error:
    \begin{minted}{text}
error[E0384]: cannot assign to immutable argument `a`
 --> src/lib.rs:2:5
  |
2 |     a = 5;
  |     ^^^^^ cannot assign to immutable argument
    \end{minted}
\end{frame}

\begin{frame}[fragile]
    \frametitle{Functions}
    \begin{minted}[linenos]{rust}
fn some_function(mut a: usize) -> usize {
    a = 5;
    a
}
    \end{minted}
    \pause
    \begin{minted}[linenos]{rust}
fn main() {
    let value = 10;
    some_function(value);
    println!("{value}");
}
    \end{minted}
    Output:
    \pause
    \begin{minted}{text}
10
    \end{minted}
    \pause
    \texttt{value} is not defined as mut, the compiler guarantees non-mutability.
\end{frame}

\begin{frame}[fragile]
    \frametitle{Tuples}
    \begin{itemize}
        \item We now look into more complex types.
        \item A tuple is a fixed-size list of heterogenoeous values.
    \end{itemize}
    \pause
    \begin{minted}[linenos]{rust}
fn main() {
    let a: (u64, char, f32) = (32, 'f', 12.47);
    let b = (13, true, 15);  // type inference works for tuples
    let c: u32 = (5);  // (5) is not a tuple
    let d: (u32,) = (5,);  // one-sized tuples look a bit odd
    let e = ();  // the unit type is just a zero-sized tuple!
}
    \end{minted}
    \pause
    Tuples are convenient to return multiple values from a function.
\end{frame}

\begin{frame}[fragile]
    \frametitle{Tuples}
    You can access fields through ``tuple indexing''.   
    \pause
    \begin{minted}[linenos]{rust}
fn main() {
    let a: (u64, char, f32) = (32, 'f', 12.47);

    let num = a.0;
    let letter = a.1;
    let float = a.2;
}
    \end{minted}
\end{frame}

\begin{frame}[fragile]
    \frametitle{Tuples}
    Alternatively, you can destructure a tuple.
    \pause
    \begin{minted}[linenos]{rust}
fn main() {
    let a: (u64, char, f32) = (32, 'f', 12.47);

    let (num, letter, float) = a;
}
    \end{minted}
\end{frame}

\begin{frame}[fragile]
    \frametitle{Advanced Printing}
    \begin{itemize}
        \item Tuples have no default formatter (we noticed already when trying to print the unit value).
        \item Tuples support debug formatting.
    \end{itemize}
    \pause
    \begin{minted}[linenos]{rust}
fn main() {
    println!("{:?}", ());
    
    let a = (32, 'f', 12.47);
    println!("{a:?}")
}
    \end{minted}
    Output:
    \begin{minted}{text}
()
(32, 'f', 12.47)
    \end{minted}
\end{frame}

\begin{frame}[fragile]
    \frametitle{Advanced Printing}
    \begin{itemize}
        \item The \texttt{:} specifies that formatting options follow.
        \item The \texttt{?} specifies debug formatting.
        \item There are a lot more formatting options (\href{https://doc.rust-lang.org/std/fmt/}{fmt documentation}).
    \end{itemize}
    \pause
    \begin{minted}[linenos]{rust}
fn main() {
    println!("{:.2}", 3.14159);  // precision of floats
    println!("{:b}", 66);  // binary
    println!("{:x}", 254);  // hexadecimal
}
    \end{minted}
    Output:
    \begin{minted}{text}
3.14
1000010
fe
    \end{minted}
\end{frame}

\begin{frame}[fragile]
    \frametitle{Arrays}
    \begin{minted}[linenos]{rust}
fn main() {
    let integer_array: [i32; 4] = [-23, 492, 8001, -3];
    let boolean_array = [true, false, false ];  // type inference works
    let tuple_array: [(u64, f64); 2] = [(0, 0.0), (1, 1.0)];
}
    \end{minted}
    \pause
    \begin{itemize}
        \item The size is part of an array's type.
        \item \texttt{[i32; 4]} is a different type than \texttt{[i32; 5]}.
        \pause
        \item The compiler needs to know the size at compile time.
    \end{itemize}
\end{frame}

\begin{frame}[fragile]
    \frametitle{Arrays}
    The \textit{repeat expression} allows easy initialization of large arrays.
    \pause
    \begin{minted}[linenos]{rust}
fn main() {
    let integer_array: [f64; 1024] = [0.0; 1024];
    let bool_array: [bool; 1000000] = [true; 1000000];
    let array_array: [[i32; 30]; 500] = [[0; 30]; 500];
}
    \end{minted}
\end{frame}

\begin{frame}[fragile]
    \frametitle{Arrays}
    \begin{itemize}
        \item The \texttt{[]} operator allows index-based access to array elements.
        \item The index must be of type \texttt{usize}.
    \end{itemize}
    \pause
    \begin{minted}[linenos]{rust}
fn main() {
    let array = ['a'; 100];

    println!("{}", array[25]);
}
    \end{minted}
    \pause
    \begin{itemize}
        \item Out-of-bounds accesses are not memory safe.
        \item The compiler checks indices either at compile time or inserts runtime checks.
    \end{itemize}
\end{frame}

\begin{frame}[fragile]
    \frametitle{Arrays}
    \begin{itemize}
        \item Arrays do not have a default formatter.
        \item Debug formatter works.
    \end{itemize}
    \pause
    \begin{minted}[linenos]{rust}
fn main() {
    let integer_array: [i32; 4] = [-23, 492, 8001, -3];

    println!("{integer_array:?}");
    println!("{} {}", [4.2, 3.96], [false; 5]);
}
    \end{minted}
    \pause
    Output:
    \begin{minted}{text}
[-23, 492, 8001, -3]
[4.2, 3.96] [false, false, false, false, false]
    \end{minted}
\end{frame}

\begin{frame}[fragile]
    \frametitle{Examples: Arrays, Tuples, Functions, \texttt{while}}
    \begin{minted}[linenos]{rust}
fn min_and_max(array: [u64; 10]) -> (u64, u64) {
    let mut min = u64::MIN;
    let mut max = u64::MAX;
    let mut i = 0;
    while i < 10 {
        let val = array[i];
        if val < min {
            min = val;
        }
        if val > max {
            max = val;
        }
        i += 1;
    }
    (min, max)
}
    \end{minted}
\end{frame}

\begin{frame}[fragile]
    \frametitle{Examples: Arrays, Tuples, Functions, \texttt{while}}
    \begin{minted}[linenos]{rust}
fn all(array: [bool; 10]) -> bool {
    let mut i = 0;
    while i < 10 {
        if !array[i] {
            return false;
        }
        i += 1;
    }
    true
}
    \end{minted}
\end{frame}

\begin{frame}[fragile]
    \frametitle{Structs}
    \begin{itemize}
        \item We are moving to custom data types now.
        \item A \texttt{struct} is a group of named or unnamed values.
    \end{itemize}
    \pause
    \begin{minted}[linenos]{rust}
struct Point(f64, f64);  // x and y coordinates
    \end{minted}
    \pause
    \begin{minted}[linenos]{rust}
struct Rectangle {
    top_left: Point,
    width: f64,
    height: f64,  // the last comma is optional but idiomatic
}
    \end{minted}
    \pause
    The variant with unnamed fields must be terminated with \texttt{;}, the other variant must not!
\end{frame}

\begin{frame}[fragile]
    \frametitle{Structs}
    Instantiating values looks very similar to the type definitions.
    \pause
    \begin{minted}[linenos]{rust}
fn main() {
    let point = Point(4., 5.);
    let rect = Rectangle {
        top_left: point,
        width: 10.,
        height: 10.,
    }
}
    \end{minted}
\end{frame}

\begin{frame}[fragile]
    \frametitle{Structs}
    When an initializer value is bound to the matching name, you can use \textit{field init shorthand} syntax.
    \pause
    \begin{minted}[linenos]{rust}
fn main() {
    let top_left = Point(4., 5.);
    let width = 10.;
    let height = 10.;
    let rect = Rectangle {
        top_left,
        width,
        height,
    }
}
    \end{minted}
\end{frame}

\begin{frame}[fragile]
    \frametitle{Structs}
    For field access, a \texttt{struct} with unnamed fields behaves like a tuple.
    \pause
    \begin{minted}[linenos]{rust}
fn print_point(point: Point) {
    println!("Point at ({}, {})", point.0, point.1);
}
    \end{minted}
    \pause
    Destructuring also works.
    \begin{minted}[linenos]{rust}
fn print_point(point: Point) {
    let Point(x, y) = point;
    println!("Point at ({x}, {y})");
}
    \end{minted}
\end{frame}

\begin{frame}[fragile]
    \frametitle{Structs}
    If a \texttt{struct} with unnamed fields behaves like a tuple but is more verbose, what's the point?
    \pause
    \begin{minted}[linenos]{rust}
struct Point(f64, f64);  // (x, y), point in 2d space
struct Direction(f64, f64);  // (x, y), direction/vector in 2d space
struct Range(f64, f64);  // (min, max), a range in 1d space
    \end{minted}
    \begin{itemize}
        \item The types store the same data.
        \pause
        \item Using \texttt{struct}s instead of tuples lets the compiler check that you do not accidently pass a \texttt{Direction} to where a \texttt{Point} is expected.
    \end{itemize}
\end{frame}

\begin{frame}[fragile]
    \frametitle{Structs}
    Access fields of a named \texttt{struct} by name.
    \pause
    \begin{minted}[linenos]{rust}
fn bottom_right(rect: Rectangle) -> Point {
    let x = rect.top_left.0 + rect.width;
    let y = rect.top_left.1 + rect.height;

    Point(x, y)
}
    \end{minted}
\end{frame}

\begin{frame}[fragile]
    \frametitle{Structs}
    You can also desctructure a \texttt{struct} with named fields.
    \pause
    \begin{minted}[linenos]{rust}
fn bottom_right(rect: Rectangle) -> Point {
    let Rectangle { top_left, width, height } = rect;

    Point(top_left.0 + width, top_left.1 + height)
}
    \end{minted}
    \pause
    \begin{minted}[linenos]{rust}
fn bottom_right(rect: Rectangle) -> Point {
    let Rectangle { top_left: p, width: w, height: h } = rect;

    Point(p.0 + w, p.1 + h)
}
    \end{minted}
\end{frame}

\begin{frame}[fragile]
    \frametitle{Structs}
    \begin{itemize}
        \item By default, \texttt{struct}s offer neither normal nor debug formatting.
        \item For now, we tag custom types with \texttt{\#[derive(Debug)]} without understanding what it does.
    \end{itemize}
    \pause
    \begin{minted}[linenos]{rust}
#[derive(Debug)]
struct Point(f64, f64);

#[derive(Debug)]
struct Rectangle {
    top_left: Point,
    width: f64,
    height: f64,
}
    \end{minted}
\end{frame}

\begin{frame}[fragile]
    \frametitle{Structs}
    \begin{minted}[linenos]{rust}
#[derive(Debug)]
struct Point(f64, f64);

fn main() {
    println("{:?}", Point(3., 4.))
}
    \end{minted}
    \pause
    Output:
    \begin{minted}{text}
Point(3.0, 4.0)
    \end{minted}
\end{frame}

\begin{frame}[fragile]
    \frametitle{Structs}
    \begin{minted}[linenos]{rust}
#[derive(Debug)]
struct Rectangle {
    <...>
}
fn main() {
    let rect = Rectangle {
        top_left: Point(3., 4.),
        width: 10.,
        height: 10.,
    };
    println!("{rect:?}");
}
    \end{minted}
    \pause
    Output:
    \begin{minted}{text}
Rectangle { top_left: Point(3.0, 4.0), width: 10.0, height: 10.0 }
    \end{minted}
\end{frame}

\end{document}
